\chapter{Introdução}

O avanço das tecnologias da informação e da comunicação tem transformado de maneira significativa os processos produtivos, administrativos e comerciais em diferentes setores da economia. A digitalização de atividades antes realizadas manualmente tornou-se um fator relevante para a organização, a competitividade e a sustentabilidade das empresas, independentemente de seu porte ou segmento de atuação.

No cenário atual, caracterizado pela necessidade de agilidade, precisão e confiabilidade na gestão empresarial, pequenos negócios enfrentam o desafio de modernizar suas rotinas operacionais sem, necessariamente, dispor de sistemas complexos ou de alto custo. A adoção de sistemas informatizados voltados ao controle de estoque e ao registro de vendas representa uma alternativa para reduzir falhas humanas, centralizar informações e oferecer maior suporte à tomada de decisão.

Inserida nesse contexto, a loja Tropical Mix, localizada no município de Guarulhos, atua na comercialização de insumos para máquinas de sorvete expresso e produtos relacionados à fabricação de sorvetes. Antes da proposta deste projeto, parte do controle de estoque e vendas era realizada de forma manual, o que poderia ocasionar inconsistências nos registros, retrabalho, dificuldade de acompanhamento das movimentações e limitações no acesso a informações atualizadas sobre produtos e vendas.

Diante dessa realidade, o presente trabalho propõe o desenvolvimento de um sistema web para gerenciamento de estoque e vendas, projetado especificamente para atender às necessidades operacionais da loja Tropical Mix. A aplicação tem como finalidade centralizar o cadastro de produtos, clientes, fornecedores e usuários, registrar vendas, acompanhar entradas e saídas de mercadorias e disponibilizar informações gerenciais para apoio à administração da empresa.

O sistema foi estruturado a partir de uma arquitetura em camadas, contemplando front-end, back-end e banco de dados. A interface foi desenvolvida em React, a camada de aplicação utiliza Python com Flask, e a persistência dos dados é realizada em MySQL. A comunicação entre as camadas ocorre por meio de uma API REST, utilizando dados em formato JSON, o que contribui para a organização, modularidade e manutenção da solução.

Assim, o projeto busca contribuir para a digitalização dos processos internos da Tropical Mix, oferecendo uma solução acessível, funcional e adequada ao contexto de pequenos negócios. Além disso, por se tratar de uma iniciativa acadêmica e extensionista, o desenvolvimento do sistema possibilita a aplicação prática de conhecimentos relacionados à engenharia de software, banco de dados, desenvolvimento web, gestão de projetos e documentação técnica.

\section{Objetivos}

O presente projeto tem como objetivo desenvolver um sistema web para gerenciamento de estoque e vendas da loja Tropical Mix, com foco na organização dos processos internos, na redução de falhas operacionais e na centralização das informações relacionadas a produtos, clientes, fornecedores, usuários e transações comerciais.

\subsection{Objetivo geral}

Desenvolver uma aplicação web para auxiliar o controle de estoque e o registro de vendas da loja Tropical Mix, substituindo processos manuais por uma solução informatizada, acessível por navegador e integrada a um banco de dados relacional.

\subsection{Objetivos específicos}

Para alcançar o objetivo geral, foram definidos os seguintes objetivos específicos:

\begin{itemize}
    \item Levantar as principais necessidades operacionais da loja Tropical Mix relacionadas ao controle de estoque e ao registro de vendas;
    
    \item Definir os requisitos funcionais e não funcionais do sistema, considerando aspectos de usabilidade, segurança, confiabilidade e manutenção;
    
    \item Projetar a arquitetura da aplicação, contemplando a separação entre front-end, back-end e banco de dados;
    
    \item Desenvolver uma interface web para autenticação de usuários, cadastro de produtos, clientes e fornecedores, controle de estoque, registro de vendas e visualização de informações gerenciais;
    
    \item Implementar a comunicação entre a interface e a API por meio de requisições HTTP e troca de dados em formato JSON;
    
    \item Estruturar o banco de dados relacional para armazenar informações de usuários, produtos, fornecedores, clientes, pedidos e movimentações de estoque;
    
    \item Validar os principais fluxos da aplicação, verificando se as funcionalidades implementadas atendem às necessidades definidas para o sistema.
\end{itemize}

\section{Problema e Solução Proposta}

Atualmente, pequenos negócios que realizam o controle de estoque e vendas de forma manual podem enfrentar dificuldades relacionadas à organização, à atualização e à confiabilidade das informações. No caso da loja Tropical Mix, localizada em Guarulhos, o acompanhamento manual das entradas e saídas de produtos pode gerar inconsistências nos registros, retrabalho, dificuldade para identificar a disponibilidade real dos itens e limitações no acompanhamento das vendas realizadas.

Esse cenário também dificulta a consulta rápida de informações importantes para a gestão, como produtos com baixa quantidade em estoque, itens próximos ao vencimento, histórico de vendas, fornecedores vinculados aos produtos e dados cadastrais de clientes. A ausência de uma ferramenta centralizada pode comprometer a eficiência operacional da loja, uma vez que as informações ficam mais sujeitas a erros humanos e a processos repetitivos.

Diante desse problema, a solução proposta consiste no desenvolvimento de um sistema web voltado ao gerenciamento interno da Tropical Mix. A aplicação tem como objetivo centralizar as principais operações relacionadas ao controle de estoque e ao registro de vendas, permitindo que usuários autorizados cadastrem produtos, clientes, fornecedores e usuários, registrem vendas, acompanhem quantidades disponíveis e consultem informações gerenciais.

O sistema foi projetado para ser acessado por navegador web, utilizando uma arquitetura composta por front-end, back-end e banco de dados. A interface permite a interação do usuário com as funcionalidades do sistema, enquanto a API realiza o processamento das regras de negócio e a comunicação com o banco de dados. Essa separação favorece a organização da aplicação, a manutenção do código e a evolução futura das funcionalidades.

A proposta busca substituir processos manuais por uma solução informatizada, simples e adequada à realidade de um pequeno negócio. Com isso, espera-se reduzir falhas operacionais, melhorar a rastreabilidade das informações, facilitar o acompanhamento do estoque e das vendas e oferecer maior suporte à tomada de decisão pela administração da loja.

\section{Justificativa}

A digitalização dos processos de gestão ainda é um desafio para grande parte das micro e pequenas empresas brasileiras. De acordo com o \textit{1º Panorama da Gestão de Despesas Corporativas no Brasil} \cite{contasimples2025}, 39\% das MPMEs ainda realizam o controle de despesas de forma manual, principalmente por questões de custo e familiaridade com métodos tradicionais. Esse cenário evidencia a necessidade de soluções tecnológicas acessíveis, simples e adaptadas à realidade de pequenos negócios.

No contexto da loja \textbf{Tropical Mix}, o desenvolvimento de um sistema web para controle de estoque e vendas justifica-se pela necessidade de centralizar informações operacionais e tornar os processos internos mais organizados, seguros e eficientes. O acompanhamento manual de produtos, entradas, saídas e vendas pode comprometer a confiabilidade dos registros, gerar retrabalho e dificultar a consulta de informações atualizadas para apoio à gestão.

Além disso, a proposta busca oferecer uma solução compatível com a realidade financeira e operacional de pequenos empreendimentos. Diferentemente de sistemas corporativos mais robustos, que podem apresentar custos elevados e maior complexidade de implantação, o sistema proposto prioriza funcionalidades essenciais, interface intuitiva e uso de tecnologias amplamente disponíveis, como React, Flask e MySQL. Dessa forma, o projeto busca promover eficiência, agilidade e confiabilidade na gestão de estoque e vendas.

Além de sua relevância técnica, o projeto possui caráter extensionista, pois aplica o conhecimento acadêmico em benefício da comunidade, fortalecendo microempresas — responsáveis por mais de 70\% dos empregos formais no Brasil \cite{sebrae2024}. Ao atender uma demanda real da loja Tropical Mix, o projeto aproxima o ambiente acadêmico das necessidades de pequenos negócios, contribuindo tanto para a formação prática dos estudantes quanto para a modernização dos processos da empresa.

Portanto, o desenvolvimento do sistema justifica-se pela combinação entre necessidade operacional, viabilidade tecnológica e impacto social. A proposta contribui para a digitalização de processos internos da Tropical Mix, ao mesmo tempo em que permite aos estudantes aplicar conhecimentos relacionados a desenvolvimento web, banco de dados, engenharia de software, gestão de projetos e documentação técnica.


\section{Análise da Concorrência}
Para compreender o posicionamento do sistema proposto no mercado, foi realizada uma análise comparativa entre soluções existentes que oferecem funcionalidades semelhantes voltadas à gestão de estoque e controle de vendas. O objetivo dessa análise é identificar os diferenciais competitivos, pontos fortes e limitações dos sistemas disponíveis, de modo a evidenciar como o sistema desenvolvido para a loja \textbf{Tropical Mix} se destaca por atender às necessidades específicas de pequenos negócios que buscam uma solução simples, eficiente e de baixo custo. 

A seguir, são apresentadas as principais características de três plataformas relevantes no segmento: Luccrei PDV, PedidoOK e WM10 ERP.
\subsection{Luccrei PDV}

A \textbf{Luccrei PDV} é uma plataforma voltada à gestão comercial e controle de vendas. Seu foco é proporcionar praticidade no atendimento ao cliente, com um sistema de frente de caixa rápido e eficiente. Dentre as principais funcionalidades, destacam-se:

\begin{itemize}
    \item \textbf{Registro de vendas} e emissão de notas fiscais;
    \item \textbf{Operações offline}, permitindo vendas mesmo sem conexão à internet;
    \item \textbf{Controle de estoque} com atualização em tempo real;
    \item \textbf{Relatórios detalhados} e gestão financeira centralizada.
\end{itemize}

A plataforma é direcionada principalmente a pequenos e médios comércios que necessitam de um sistema ágil e intuitivo. No entanto, apresenta limitações no que diz respeito à personalização e integração com outros sistemas de análise de dados.

\subsection{PedidoOK}

O \textbf{PedidoOK} é um sistema completo para controle de vendas, estoque e emissão de documentos fiscais (NF-e, NFC-e e boletos). O sistema é voltado tanto para \textit{micro, pequenas e médias empresas} quanto para distribuidores e representantes comerciais. Suas principais características incluem:

\begin{itemize}
    \item \textbf{Gestão integrada de vendas, estoque e finanças};
    \item \textbf{Aplicativo Android} para controle de pedidos e catálogo de produtos;
    \item \textbf{Integração com marketplaces}, bancos e meios de pagamento;
    \item Emissão automatizada de documentos fiscais.
\end{itemize}

O diferencial do PedidoOK está na mobilidade e na integração com plataformas externas. Entretanto, a necessidade de assinatura mensal e a curva de aprendizado podem representar barreiras para pequenos negócios que buscam soluções mais simples e econômicas.

\subsection{WM10 ERP}

A \textbf{WM10 ERP} oferece uma solução robusta e especializada para centralizar operações de \textbf{produção, loja física, e-commerce e equipe externa}. O sistema se destaca pela sua abordagem completa, sendo direcionado a empresas de médio e grande porte que necessitam de alto nível de automação. Suas principais funcionalidades são:

\begin{itemize}
    \item \textbf{Controle de estoque e produção em tempo real};
    \item \textbf{Integração omnichannel}, unificando loja física e virtual;
    \item \textbf{Automação de processos empresariais};
    \item \textbf{Relatórios gerenciais} e dashboards de desempenho.
\end{itemize}

Apesar de ser uma plataforma avançada, o WM10 apresenta custos mais elevados e complexidade de implantação, tornando-se menos acessível para pequenos negócios locais.

\subsection{Comparativo}

A Tabela~\ref{tab:comparativo-concorrencia} apresenta um comparativo geral entre as três soluções analisadas, destacando as principais funcionalidades oferecidas pelos sistemas Luccrei PDV, PedidoOK, WM10 ERP e o sistema proposto para a loja Tropical Mix.

\begin{table}[H]
\centering
\caption{Comparativo de funcionalidades entre Tropical Mix e sistemas concorrentes.}
\label{tab:comparativo-concorrencia}
\resizebox{\textwidth}{!}{%
\begin{tabular}{|l|c|c|c|c|}
\hline
\textbf{Funcionalidade} & \textbf{Luccrei PDV} & \textbf{PedidoOK} & \textbf{WM10 ERP} & \textbf{Tropical Mix} \\ \hline
Controle de estoque & X & X & X & X \\ \hline
Emissão de notas fiscais & X & X & X &  \\ \hline
Aplicativo móvel &  & X & X &  \\ \hline
Acesso via navegador & X & X & X & X \\ \hline
Relatórios gerenciais & X & X & X & X \\ \hline
Dashboard de controle &  & X & X &  \\ \hline
Integração com plataformas externas & X & X & X &  \\ \hline
Gestão financeira & X & X & X &  \\ \hline
Gestão de clientes e vendas & X & X & X & X \\ \hline
Alertas automáticos de estoque &  &  & X & X \\ \hline
Suporte técnico dedicado &  & X & X &  \\ \hline
Custo estimado & Médio & Médio & Alto & Baixo \\ \hline
\end{tabular}%
}
\fonte{Elaborado pelos autores.}
\end{table}

A partir dessa análise, percebe-se que o sistema proposto para a loja Tropical Mix busca um equilíbrio entre simplicidade e eficiência, atendendo às necessidades específicas de um negócio de pequeno porte, sem os custos e complexidades dos sistemas corporativos mais robustos.